\begin{table}[t]
\centering
\caption{Attack outcomes by protection policy, on the same credit register and the same adversary. The release and pooled columns give the share of targets the maximum-posterior linkage names correctly, for one release and pooled over the 6 vintages. Precision is measured over the attempts the adversary acts on, those whose posterior passes 0.5, and the column headed acts is how often it acts at all. The last two columns use the same success event as the release column for the size class with the highest rate, with that class's exact 95 per cent upper bound; a release passes the subgroup control when every class's bound meets the tolerance, which needs no correction for the number of classes. Utility is the negative weighted loss on the four analyses the register is collected for. Noisy-count releases are not compared: their targets are key classes, not loans.}
\label{tab:eval}
\small
\begin{tabular}{lrrrrrrrr}
\toprule
Policy & Kept & Utility & Release & Pooled & Precision & Acts & Worst class & Upper \\
\midrule
no suppression & 27{,}167 & -0.000 & 0.089 & 0.260 & 0.716 & 0.020 & 0.248 & 0.297 \\
three contributors & 27{,}068 & -0.019 & 0.084 & 0.255 & 0.822 & 0.018 & 0.180 & 0.228 \\
complementary suppression & 27{,}027 & -0.028 & 0.080 & 0.259 & 0.746 & 0.015 & 0.156 & 0.203 \\
$k$-anonymity, $k=5$ & 26{,}885 & -0.050 & 0.074 & 0.269 & 0.771 & 0.012 & 0.183 & 0.234 \\
$k$-anonymity, $k=10$ & 26{,}019 & -0.139 & 0.070 & 0.244 & 0.800 & 0.009 & 0.141 & 0.193 \\
\bottomrule
\end{tabular}
\end{table}
